\chapter{Pipeline Source Code Structure}
\label{annex:pipeline}

This annex describes the organization of the source code repository
implementing the data pipeline described in \Cref{ch:methodology}. All
code is written in Python, using \texttt{pandas} and \texttt{numpy} for
data manipulation, \texttt{yfinance} for market data retrieval,
\texttt{Hugging Face transformers} and \texttt{torch} for FinBERT
inference, and \texttt{scikit-learn} / \texttt{hmmlearn} for clustering.
The repository is organized into four top-level packages under
\texttt{src/}, corresponding to the four major stages of the pipeline:
\texttt{news}, \texttt{market}, \texttt{sentiment} and \texttt{clustering},
plus an \texttt{eda} package containing exploratory analyses that informed
several methodological decisions (e.g., the choice of the MA21 smoothing
window in \Cref{sec:meth-smoothing}).

\section{Repository Structure}
\label{annex:pipeline-structure}

\begin{verbatim}
src/
├── news/
│   └── load_fnspid.py            Stream and cache the raw FNSPID corpus
├── market/
│   ├── download_sp500.py         Download S&P 500 OHLC prices (yfinance)
│   ├── sp500_historical.py        Historical S&P 500 constituents lookup
│   └── technical_indicators.py    Compute the 9 technical indicators
├── sentiment/
│   ├── sentiment_finbert.py       FinBERT batched sentiment scoring
│   ├── aggregate_sentiment.py     (date, ticker)-level aggregation
│   ├── aggregate_by_sector.py     (date, sector)-level aggregation + MA21
│   ├── cross_market_sentiment.py  Merge sentiment with market returns
│   └── build_clustering_dataset.py  Final clustering dataset assembly
├── eda/
│   ├── eda_gics_mapping.py         Build ticker -> GICS sector mapping
│   ├── eda_filter_sp500_news.py    Filter FNSPID news by S&P 500 membership
│   ├── eda_ticker_quality.py       Per-ticker news coverage diagnostics
│   ├── eda_news.py                 News volume / distribution analyses
│   ├── eda_market.py               Market data quality checks
│   ├── eda_news_vs_market.py       News vs. market exploratory analysis
│   ├── eda_fnspid_full_stream_analysis.py  Full-corpus FNSPID statistics
│   └── eda_sentiment_smoothing.py  MA5/MA10/MA21 noise-reduction analysis
└── clustering/
    ├── clustering.py               K-Means Models A/B/C + Gaussian HMM
    ├── silhouette_pca.py           PCA-common-space silhouette comparison
    ├── feature_importance.py       ANOVA F-statistic feature importance
    ├── plot_clusters.py             PCA, correlation and regime plots
    ├── economic_validation.py       Market-level economic metrics
    └── sector_regime_analysis.py    Sector-level economic metrics & heatmaps
\end{verbatim}

Each module follows the same general convention: input and output paths
are defined as module-level constants (typically under \texttt{data/raw/},
\texttt{data/processed/} or \texttt{outputs/}), and a \texttt{main()}
function executes the corresponding processing step, printing progress
information (row counts, date ranges, intermediate statistics) to allow
each stage to be run and verified independently. The following sections
describe each stage in more detail.

\section{Data Acquisition}
\label{annex:pipeline-acquisition}

\textbf{\texttt{news/load\_fnspid.py}} streams the FNSPID dataset
(\texttt{Zihan1004/FNSPID} on the Hugging Face Hub, corresponding to
\cite{zhang2023fnspid}) using the \texttt{datasets} library in streaming
mode, and writes the relevant columns (publication date, headline title,
stock ticker, publisher and URL) to a local CSV cache under
\texttt{data/raw/news/}. Streaming mode is used because the full corpus
contains several million rows and does not fit comfortably in memory.

\textbf{\texttt{market/download\_sp500.py}} downloads daily OHLC
(open/high/low/close/volume) data for the S\&P 500 index from Yahoo
Finance via \texttt{yfinance}, with a retry mechanism (up to 3 attempts per
ticker) to handle transient network failures, and caches the result under
\texttt{data/raw/market/}.

\textbf{\texttt{market/sp500\_historical.py}} provides a lookup function,
\texttt{sp500\_tickers\_on(date)}, which, given a \texttt{data/raw/market/sp500\_historical\_components.csv}
file containing the historical list of S\&P 500 constituents at successive
dates, returns the list of tickers that were constituents of the index on
(or most recently before) a given date. This function underpins the
membership-filtering step described in \Cref{sec:meth-sp500}.

\section{News Filtering and GICS Sector Mapping}
\label{annex:pipeline-filtering}

\textbf{\texttt{eda/eda\_gics\_mapping.py}} constructs a
ticker-to-GICS-sector mapping (\texttt{data/raw/market/ticker\_gics\_mapping.csv})
by combining a base table of known sectors with sector information
retrieved from Yahoo Finance for any remaining tickers, and normalizes
Yahoo Finance's sector names to the standard GICS sector names used
throughout this thesis (e.g., \texttt{"Healthcare"} $\to$ \texttt{"Health
Care"}, \texttt{"Financials"} $\to$ \texttt{"Financial Services"},
\texttt{"Consumer Cyclical"} $\to$ \texttt{"Consumer Discretionary"},
\texttt{"Consumer Defensive"} $\to$ \texttt{"Consumer Staples"},
\texttt{"Basic Materials"} $\to$ \texttt{"Materials"}, \texttt{"Technology"}
$\to$ \texttt{"Information Technology"}), producing the 11-sector
distribution summarized in \Cref{tab:gics-distribution}.

\textbf{\texttt{eda/eda\_filter\_sp500\_news.py}} implements the
membership-filtering step of \Cref{sec:meth-filtering}: it expands the
historical S\&P 500 constituents table into a long (date, ticker) table,
and retains only those FNSPID news rows whose (date, ticker) pair appears
in this table --- i.e., news items associated with a company that was an
S\&P 500 constituent on the date the news was published. The filtered
result is written to \texttt{data/processed/news/fnspid\_sp500.parquet}.

\section{FinBERT Sentiment Scoring}
\label{annex:pipeline-finbert}

\textbf{\texttt{sentiment/sentiment\_finbert.py}} implements the sentiment
scoring procedure of \Cref{sec:meth-finbert}. After loading the filtered
news parquet, normalizing column names (\texttt{date}, \texttt{ticker},
\texttt{title}) and dropping unused summary columns, the script loads the
\texttt{ProsusAI/finbert} model via the \texttt{transformers.pipeline} API
(using GPU acceleration when available, \texttt{torch.cuda.is\_available()}),
and processes all headlines in batches of \texttt{BATCH\_SIZE = 64}. Each
headline is truncated to 512 characters before tokenization (FinBERT's
maximum input length). For each headline, the pipeline returns class
probabilities $p_{\text{pos}}, p_{\text{neg}}, p_{\text{neu}}$, from which
the continuous sentiment score $s = p_{\text{pos}} - p_{\text{neg}}$
(\Cref{eq:sentiment-score}) and the categorical label
(\Cref{eq:sentiment-label}) are derived. The module also supports a
\texttt{DEBUG\_MODE} flag that restricts processing to a random subset of
50,000 headlines (\texttt{random\_state=42}) for fast iteration during
development; all results reported in this thesis use the full corpus
(\texttt{DEBUG\_MODE = False}).

\section{Sector-Level Aggregation and Smoothing}
\label{annex:pipeline-aggregation}

\textbf{\texttt{sentiment/aggregate\_sentiment.py}} performs an
intermediate (date, ticker)-level aggregation, computing, for each
ticker and day, the mean and standard deviation of the sentiment score
across all headlines published that day (\texttt{sentiment\_score\_mean},
\texttt{sentiment\_score\_std}), the total number of news items, and the
counts of positive/negative/neutral headlines.

\textbf{\texttt{sentiment/aggregate\_by\_sector.py}} performs the
sector-level aggregation described in \Cref{sec:meth-aggregation}: it
merges the per-ticker sentiment scores with the GICS sector mapping
(an inner join on \texttt{ticker}), groups by \texttt{(date, sector)}, and
computes \texttt{sentiment\_avg} (mean) and \texttt{sentiment\_std}
(standard deviation, \Cref{eq:sector-mean,eq:sector-std}) of the sentiment
score across all news items published about companies in that sector on
that day. Standard deviations for (date, sector) cells with a single news
item (undefined sample standard deviation) are filled with zero. The
resulting long table is then pivoted into a wide format with one column
per (statistic, sector) combination, and a rolling mean with
\texttt{SMOOTH\_WINDOW = 21} sessions (\Cref{eq:ma-smooth}) is applied to
each \texttt{sentiment\_avg} column to produce the corresponding
\texttt{\_smooth} columns used by Model C.

\textbf{\texttt{sentiment/cross\_market\_sentiment.py}} is an exploratory
module that merges the aggregated per-ticker sentiment series with daily
S\&P 500 returns and rolling volatility, used during the EDA phase to
visualize the relationship between sentiment and market dynamics
(\Cref{sec:soa-gap}); it is not part of the critical path that produces the
final clustering dataset.

\section{Technical Indicators}
\label{annex:pipeline-technical}

\textbf{\texttt{market/technical\_indicators.py}} computes the 9 technical
indicators of \Cref{sec:meth-technical} from the daily S\&P 500 OHLC
series. Log returns are computed as
$r_t = \ln(\text{close}_t / \text{close}_{t-1})$. The Relative Strength
Index is computed via \texttt{compute\_rsi()}, using exponentially-weighted
moving averages of gains and losses with a center-of-mass parameter
\texttt{com = window - 1} for a 14-session window
(\Cref{eq:rsi}). The 50- and 200-session moving-average ratios
(\texttt{MA50\_ratio}, \texttt{MA200\_ratio}, \Cref{eq:ma-ratio}), the
21- and 126-session momentum measures (\texttt{mom21}, \texttt{mom126},
\Cref{eq:momentum}), and the 30- and 252-session annualized realized
volatilities (\texttt{vol\_30d}, \texttt{vol\_252d}, scaled by
$\sqrt{252}$, \Cref{eq:realized-vol}) and rolling skewness measures
(\texttt{skew\_30d}, \texttt{skew\_252d}, \Cref{eq:skewness}) are computed
using \texttt{pandas} rolling-window operations on the log-return series.

\section{Clustering Dataset Construction}
\label{annex:pipeline-dataset}

\textbf{\texttt{sentiment/build\_clustering\_dataset.py}} performs the
final assembly step of \Cref{sec:meth-dataset}. It loads the sector-level
sentiment table and the technical indicator table, renames the technical
columns with the \texttt{\_sp500} suffix (e.g., \texttt{RSI} $\to$
\texttt{RSI\_sp500}), drops auxiliary columns (\texttt{ret}, \texttt{close})
that are not used as clustering features, and performs an inner merge on
\texttt{date}. A date cutoff (\texttt{DATE\_CUTOFF}) is applied to exclude
incomplete trailing observations, and rows with missing values --- arising
from the warm-up periods required by the longest rolling windows
(252 sessions for \texttt{vol252d\_sp500} and \texttt{skew252d\_sp500}) ---
are dropped, yielding the final 1{,}132-observation, 46-column dataset
written to \texttt{data/processed/clustering\_dataset.parquet}.

\section{Clustering, Internal Validation and Economic Validation}
\label{annex:pipeline-clustering}

\textbf{\texttt{clustering/clustering.py}} is the central module of the
clustering stage. It defines the three feature sets \texttt{FEATURES\_TECH}
(9 technical features, Model A), \texttt{FEATURES\_ALL} (\texttt{FEATURES\_TECH}
plus the 22 raw sentiment columns, Model B) and the 11 MA21-smoothed
sentiment columns (Model C), standardizes each feature set with
\texttt{StandardScaler}, and fits \texttt{sklearn.cluster.KMeans} with
$K=3$ for each of the three configurations (\Cref{eq:soa-kmeans-objective}),
storing the resulting labels as \texttt{kmeans\_tech}, \texttt{kmeans\_tech\_sent}
and \texttt{kmeans\_sent\_smooth} columns in
\texttt{outputs/clustering/clustered\_data.parquet}. The same module also
implements \texttt{run\_hmm()}, which fits
\texttt{hmmlearn.hmm.GaussianHMM(n\_components=3, covariance\_type="full",
n\_iter=1000, random\_state=42, min\_covar=1e-3)} on the standardized
\texttt{FEATURES\_TECH} and \texttt{FEATURES\_ALL} feature sets, producing
the \texttt{hmm\_tech} and \texttt{hmm\_tech\_sent} state assignments via
\texttt{model.predict()}, and saves a diagnostic plot of 30-day realized
volatility colored by HMM state
(\Cref{fig:hmm-states-modelA,fig:hmm-states-modelB}).

\textbf{\texttt{clustering/silhouette\_pca.py}} implements the
common-PCA-space silhouette comparison of
\Cref{sec:results-silhouette}: it standardizes the union of the 9 technical
and 11 smoothed sentiment features (20 features), fits a PCA retaining
components up to 80\% explained variance (9 components, 82.0\% actual
variance), projects each of the three K-Means partitions into this common
space, and computes the silhouette coefficient
(\Cref{eq:soa-silhouette}) of each partition in the projected space, saving
\texttt{fig\_silhouette\_pca\_comparison.png} and \texttt{fig\_silhouette\_scree.png}.

\textbf{\texttt{clustering/feature\_importance.py}} computes, for the
Model B partition (\texttt{kmeans\_tech\_sent}), the ANOVA F-statistic of
each of the 31 input features with respect to the 3-cluster grouping
using \texttt{sklearn.feature\_selection.f\_classif}, ranks the features by
F-statistic, and saves the results to
\texttt{outputs/clustering/validation/feature\_importance.csv} and
\texttt{fig\_feature\_importance.png} (\Cref{tab:feature-importance,fig:feature-importance}).

\textbf{\texttt{clustering/plot\_clusters.py}} centralizes the visual
diagnostics shared across models: it defines the
\texttt{CLUSTER\_COLORS} mapping ($\{0:\texttt{\#0D7A3E}$ (Bull, green),
$1:\texttt{\#C0392B}$ (Bear, red), $2:\texttt{\#E67E22}$ (Crisis, orange)$\}$)
and \texttt{CLUSTER\_NAMES} mapping used consistently across all figures in
this thesis, and implements \texttt{plot\_pca\_clusters()} (PCA scatter +
loadings + scree plot, \Cref{fig:pca-modelA,fig:pca-modelB}),
\texttt{plot\_sp500\_regimes()} (S\&P 500 price series with regime-shaded
background, \Cref{fig:sp500-regimes-modelA,fig:sp500-regimes-modelB}) and
\texttt{plot\_correlation\_matrix()} (feature correlation heatmaps,
\Cref{fig:correlation-modelA,fig:correlation-modelB}).

\textbf{\texttt{clustering/economic\_validation.py}} implements the
market-level economic validation framework of \Cref{sec:meth-validation}:
for each clustering configuration (Models A, B, C and both HMM variants)
and each regime/state label, it computes the mean daily return, the
annualized return (\Cref{eq:ann-return}), the annualized volatility
(\Cref{eq:ann-vol}), the Sharpe ratio with $r_f = 2\%$
(\Cref{eq:sharpe}), the percentage of positive days, and the best and worst
single-day returns, on the \emph{aggregate} S\&P 500 return series,
writing one CSV file per configuration to
\texttt{outputs/clustering/validation/} (e.g.,
\texttt{metrics\_kmeans\_tech\_sent.csv}, \texttt{metrics\_hmm\_tech.csv}),
summarized in \Cref{tab:market-kmeans,tab:market-hmm}.

\textbf{\texttt{clustering/sector\_regime\_analysis.py}} extends the
economic validation to the sector level (\Cref{sec:results-econ-sector}).
For each of the three K-Means models, it defines a \texttt{MODELS}
configuration dictionary and a \texttt{REGIME\_MAP = \{0: "Bull", 1:
"Bear", 2: "Crisis"\}} label mapping, reindexes the 1{,}132-observation
cluster labels onto the daily calendar of each of the 11 GICS sector return
series via forward-fill (yielding 3{,}695 sector-days per model,
\Cref{sec:results-econ-sector}), computes the same set of economic metrics
plus the maximum drawdown over (non-contiguous) regime days
(\Cref{eq:max-dd}) for every (sector, regime) combination, and writes the
results to \texttt{outputs/clustering/validation/<model>/sector\_regime\_metrics.csv}
(reproduced in full in \Cref{annex:tables}). The function
\texttt{generate\_heatmaps()} then produces, for each model, seven heatmap
figures (Sharpe ratio, maximum drawdown, percentage of positive days, best
day, worst day, annualized return and annualized volatility, by sector and
regime), of which the Sharpe ratio, maximum drawdown, best-day and
worst-day heatmaps for Model B and the Sharpe ratio heatmaps for Models A
and C are reproduced in \Cref{sec:results-econ-sector,sec:results-modelC}.
